\section{Выводы}В ходе выполнения лабораторной работы по курсу «Дискретный анализ» была реализована система эффективного поиска паттерна в текстовых данных на основе \textbf{алгоритма Кнута-Морриса-Пратта}.
Основные результаты работы:
\begin{itemize}
    \item Реализован препроцессинг искомого шаблона с использованием \textbf{алгоритма Гусфилда} для построения $Z$-блоков, на основе которых был сформирован оптимизированный массив сдвигов (массив сильных префиксов).
    \item Спроектирована архитектура \textbf{потокового поиска} (stream processing), позволяющая обрабатывать неограниченные объёмы входных данных без их полной загрузки в оперативную память, что обеспечивает пространственную сложность $O(n)$, где $n$ - количество слов в паттерне.
    \item Внедрён механизм \textbf{кольцевого буфера} для хранения метаданных (номера строк и позиций слов), что позволило эффективно восстанавливать координаты начала вхождения паттерна при его полном сопоставлении.
\end{itemize}